% Options for packages loaded elsewhere
\PassOptionsToPackage{unicode}{hyperref}
\PassOptionsToPackage{hyphens}{url}
%
\documentclass[
]{article}
\usepackage{amsmath,amssymb}
\usepackage{iftex}
\ifPDFTeX
  \usepackage[T1]{fontenc}
  \usepackage[utf8]{inputenc}
  \usepackage{textcomp} % provide euro and other symbols
\else % if luatex or xetex
  \usepackage{unicode-math} % this also loads fontspec
  \defaultfontfeatures{Scale=MatchLowercase}
  \defaultfontfeatures[\rmfamily]{Ligatures=TeX,Scale=1}
\fi
\usepackage{lmodern}
\ifPDFTeX\else
  % xetex/luatex font selection
\fi
% Use upquote if available, for straight quotes in verbatim environments
\IfFileExists{upquote.sty}{\usepackage{upquote}}{}
\IfFileExists{microtype.sty}{% use microtype if available
  \usepackage[]{microtype}
  \UseMicrotypeSet[protrusion]{basicmath} % disable protrusion for tt fonts
}{}
\makeatletter
\@ifundefined{KOMAClassName}{% if non-KOMA class
  \IfFileExists{parskip.sty}{%
    \usepackage{parskip}
  }{% else
    \setlength{\parindent}{0pt}
    \setlength{\parskip}{6pt plus 2pt minus 1pt}}
}{% if KOMA class
  \KOMAoptions{parskip=half}}
\makeatother
\usepackage{xcolor}
\usepackage[margin=1in]{geometry}
\usepackage{graphicx}
\makeatletter
\newsavebox\pandoc@box
\newcommand*\pandocbounded[1]{% scales image to fit in text height/width
  \sbox\pandoc@box{#1}%
  \Gscale@div\@tempa{\textheight}{\dimexpr\ht\pandoc@box+\dp\pandoc@box\relax}%
  \Gscale@div\@tempb{\linewidth}{\wd\pandoc@box}%
  \ifdim\@tempb\p@<\@tempa\p@\let\@tempa\@tempb\fi% select the smaller of both
  \ifdim\@tempa\p@<\p@\scalebox{\@tempa}{\usebox\pandoc@box}%
  \else\usebox{\pandoc@box}%
  \fi%
}
% Set default figure placement to htbp
\def\fps@figure{htbp}
\makeatother
\setlength{\emergencystretch}{3em} % prevent overfull lines
\providecommand{\tightlist}{%
  \setlength{\itemsep}{0pt}\setlength{\parskip}{0pt}}
\setcounter{secnumdepth}{-\maxdimen} % remove section numbering
\usepackage{bookmark}
\IfFileExists{xurl.sty}{\usepackage{xurl}}{} % add URL line breaks if available
\urlstyle{same}
\hypersetup{
  pdftitle={Contact},
  hidelinks,
  pdfcreator={LaTeX via pandoc}}

\title{Contact}
\author{}
\date{\vspace{-2.5em}}

\begin{document}
\maketitle

\pandocbounded{\includegraphics[keepaspectratio]{~/Documents/GitHub/kunxu.github.io/images/meandrobot.png}}

Department of Media Production, Management, and Technology\\
College of Journalism and Communications\\
University of Florida\\
1885 Stadium Rd, Gainesville, FL 32611\\
Office: Weimer Hall \#3065\\
Office Phone: 352-392-0435

Visiting Scholar Office 432 Tech Impact and Policy Center Stanford
University 616 Jane Stanford Way, Stanford, CA 94305

Email: kun {[}dot{]} xu {[}at{]} ufl {[}dot{]} edu

中文介绍：

本人现为美国佛罗里达大学新闻与传播学院副教授(终身教职)，兼任学院\href{https://www.jou.ufl.edu/metl/}{媒介效果与技术实验室}主任以及\href{https://www.jou.ufl.edu/researchlab/}{学院实验室}主任。2026年9月起在斯坦福大学\href{https://tip.fsi.stanford.edu/}{技术影响力与政策中心}以及\href{https://sml.stanford.edu/}{社交媒体实验室}做访问学者。研究领域聚焦于社交机器人(social
robots)，虚拟现实(VR)和增强现实(AR)。特别关注有关交互临场感 (social
presence) 与空间临场感 (spatial presence)
的理论建设与实证研究。教学领域包括人机传播(human-machine communication),
人机互动(human-computer interaction), 计算机辅助传播(computer-mediated
communication), 用户体验(user experience)。

\end{document}
